\section{结论}

角分辨傅里叶算子超表面设计更好地被表述为响应空间逆向问题，而非一系列不连贯的任务优化。这是本文的中心结论。通过将目标视为角度、波长和偏振上的耦合场，本研究将二阶微分、高阶微分、角度滤波和偏振相关操作重新框定为一种物理规范语言的不同实例。相关的逆向设计策略自然地是生成式的，因为可行集是多模态的，且它必然是物理一致的，因为仅替代一致性不能证明算子行为。

在该表述内，当前结果支持由以下部分组成的设计循环：仿真基础监督、替代辅助候选集中、基于条件扩散的逆向生成、以及基于RCWA的验证和重排序。因此，科学贡献不是单独器件的目录。它是自由形态傅里叶算子超表面的响应中心逆向设计范式，其中目标规范、候选生成和最终选择保持与器件的验证散射行为相关联。从该角度看，未来进展可能更少依赖于累积孤立的算子演示，而更多依赖于学习、采样和验证结构化的光学响应流形。
